\documentclass[tikz,border=12pt]{standalone}
\usepackage[UTF8,fontset=fandol]{ctex}
\usepackage{amsmath,amssymb}
\usetikzlibrary{arrows.meta}
\definecolor{teal}{HTML}{4F8FA5}
\definecolor{orange}{HTML}{EE995B}
\definecolor{coral}{HTML}{C95B5B}
\definecolor{violet}{HTML}{8A74B5}
\begin{document}
\begin{tikzpicture}[x=1cm,y=1cm,>=Latex, every node/.style={font=\small},flow/.style={->,thick,black!55}]


\node[font=\large] at (8,2.9) {$H=\phi(X[W_1^{(0)}|W_1^{(1)}])=[H^{(0)}|H^{(1)}]$};
\node[font=\large] at (8,1.15) {$Y=[H^{(0)}|H^{(1)}]\begin{bmatrix}W_2^{(0)}\\W_2^{(1)}\end{bmatrix}=P^{(0)}+P^{(1)},\quad P^{(r)}=H^{(r)}W_2^{(r)}$};
\node[font=\small\bfseries,text=black!60] at (8,0) {SP 区域：同一 TP 组内，每卡负责一半 token，保留完整隐藏维};
\filldraw[fill=teal!60,draw=black!60] (1,-2.3) rectangle ++(1.6,1);
\node at (1.8,-2.6999999999999997) {$X_0$};
\node[font=\scriptsize,text=black!60] at (1.8,-3.1999999999999997) {$(L/2)\times d$};
\node at (1.8,-1.7999999999999998) {Rank 0};
\filldraw[fill=teal!60,draw=black!60] (4,-2.3) rectangle ++(1.6,1);
\node at (4.8,-2.6999999999999997) {$X_1$};
\node[font=\scriptsize,text=black!60] at (4.8,-3.1999999999999997) {$(L/2)\times d$};
\node at (4.8,-1.7999999999999998) {Rank 1};

\node[draw=black!50,minimum width=3cm,minimum height=1cm,align=center] (ag) at (9,-1.8) {AllGather\\沿序列轴拼接};
\draw[flow] (6.1,-1.8) -- (ag.west);
\node[align=left,anchor=west] at (11.2,-1.8) {每卡都获得相同的 $X$\\$X=\operatorname{concat}_{\mathrm{seq}}(X_0,X_1)$};
\draw[flow] (ag.south) -- (9,-3.7) -- (-.6,-3.7) -- (-.6,-6.1) -- (.5,-6.1);
\node[font=\small\bfseries,text=black!60] at (8,-4.5) {TP 区域：两卡执行同一条路径，但使用各自的权重分片（$r=0,1$）};
\filldraw[fill=teal!60,draw=black!60] (0.5,-7.1) rectangle ++(1.6,2);
\node at (1.3,-7.5) {$X$};
\node[font=\scriptsize,text=black!60] at (1.3,-8.0) {$L\times d$};

\filldraw[fill=orange!60,draw=black!60] (6.9,-7.1) rectangle ++(3.2,2);
\node at (8.5,-7.5) {$H^{(r)}$};
\node[font=\scriptsize,text=black!60] at (8.5,-8.0) {$L\times(f/2)$};

\filldraw[fill=coral!60,draw=black!60] (15,-7.1) rectangle ++(1.6,2);
\node at (15.8,-7.5) {$P^{(r)}$};
\node[font=\scriptsize,text=black!60] at (15.8,-8.0) {$L\times d$};


\node[draw=black!50,minimum width=3.2cm,minimum height=1cm,align=center] (col) at (4.5,-6.1) {列并行 Linear + $\phi$\\$H^{(r)}=\phi(XW_1^{(r)})$};
\node[font=\scriptsize,text=black!60] at (4.5,-7.65) {$W_1^{(r)}:\ d\times(f/2)$};
\node[draw=black!50,minimum width=3.2cm,minimum height=1cm,align=center] (row) at (12.55,-6.1) {行并行 Linear\\$P^{(r)}=H^{(r)}W_2^{(r)}$};
\node[font=\scriptsize,text=black!60] at (12.55,-7.65) {$W_2^{(r)}:\ (f/2)\times d$};
\draw[flow] (2.1,-6.1) -- (col.west);
\draw[flow] (col.east) -- (6.9,-6.1);
\draw[flow] (10.1,-6.1) -- (row.west);
\draw[flow] (row.east) -- (15,-6.1);
\draw[flow] (16.6,-6.1) -- (17.3,-6.1) -- (17.3,-10.4) -- (8.5,-10.4) -- (8.5,-11.2);
\node[font=\small\bfseries,text=black!60,fill=white] at (8,-9.75) {两张卡的部分和送入 ReduceScatter，回到 SP 区域};
\filldraw[fill=coral!60,draw=black!60] (0.5,-12.7) rectangle ++(1.6,2);
\node at (1.3,-13.1) {$P^{(0)}$};
\node[font=\scriptsize,text=black!60] at (1.3,-13.6) {$L\times d$};

\filldraw[fill=coral!60,draw=black!60] (3.4,-12.7) rectangle ++(1.6,2);
\node at (4.2,-13.1) {$P^{(1)}$};
\node[font=\scriptsize,text=black!60] at (4.2,-13.6) {$L\times d$};

\filldraw[fill=violet!60,draw=black!60] (11.7,-12.2) rectangle ++(1.6,1);
\node at (12.5,-12.6) {$Y_0$};
\node[font=\scriptsize,text=black!60] at (12.5,-13.1) {$(L/2)\times d$};
\node at (12.5,-11.7) {Rank 0};
\filldraw[fill=violet!60,draw=black!60] (14.7,-12.2) rectangle ++(1.6,1);
\node at (15.5,-12.6) {$Y_1$};
\node[font=\scriptsize,text=black!60] at (15.5,-13.1) {$(L/2)\times d$};
\node at (15.5,-11.7) {Rank 1};

\node[font=\Large] at (2.75,-11.7) {$+$};
\node[draw=black!50,minimum width=4cm,minimum height=1cm,align=center] (rs) at (8.5,-11.7) {ReduceScatter\\跨卡求和，再按序列切分};
\draw[flow] (5.45,-11.7) -- (rs.west);
\draw[flow] (rs.east) -- (11.2,-11.7);
\node[text=black!60] at (8,-14.25) {$Y_i=(P^{(0)}+P^{(1)})[\mathcal T_i,:]$：接下来本地做 Dropout、残差与下一次 LayerNorm};
\node[fill=black!4,anchor=north west,text width=16.4cm,inner sep=10pt,align=left] at (-.2,-15.1) {
\textbf{维度}\quad $L$：序列长度；$d$：隐藏维；$f$：MLP 中间维。省略 batch、bias；$\phi$ 为逐元素激活。\\[5pt]
\textbf{对象}\quad $X_i$ 是 LayerNorm 后的序列片段；$H^{(r)}$ 是通道片段；$P^{(r)}$ 是实值部分和，并非 $Y$ 的切片。\\[5pt]
\textbf{机制}\quad SP 与 TP 复用同两张卡。LayerNorm 只需单个 token 的完整 $d$；TP 矩阵乘法需要各卡的权重分片。};
\end{tikzpicture}
\end{document}
